% Cote 119: offprints, reports, courses, CV and letters (typed and printed part): archive cote 119, site page 199 (PDF page 200).
% typescript. Draft compared: none.
% AI-transcribed by Claude Opus 5.5 (claude-opus-5-5, adaptive thinking, effort high), Message Batches API, batch r119_full (msgbatch_01DwZNm5Y7YsKBWyAN4sFs8w), 2026-11-13; not yet checked by eye.
\scanpage{119/199}
\mspage{- 4 -}

de la \uline{descente} (développé de façon systématique par Giraud [*]).

\uline{Références} : J. Giraud, P. Gabriel.

b) \uline{Algèbre commutative}: Dans le langage géométrique des ``schémas'', l'algèbre
commutative peut être considérée comme étant, essentiellement, l'étude locale
des schémas. C'est ainsi que [15],et notamment le Chap. IV de cet ouvrage, con-%
tient de très nombreux résultats nouveaux d'algèbre commutative, dont il ne
peut être question ici d'énumérer même les plus couramment utilisés. Notons
seulement ici, en algèbre locale, la notion d'anneau \uline{excellent} et ses proprié-%
tés de permanence (dont l'absence constituait sans doute la lacune la plus
marquante de l'ouvrage de M. Nagata sur les anneaux locaux).

\uline{Références} : M. Nagata, P. Samuel, M. Raynaud, O. Zariski.

c) \uline{Théorie du groupe de Brauer}: Mes contributions découlent pour l'essentiel
de l'application de la cohomologie étale (développée dans [16, SGA 3]) à la
théorie du groupe de Brauer. J'ai fait un exposé d'ensemble sur les résultats
connus sur ce groupe dans [18].

\uline{Références} : M. Artin, J. Tate, J.P. Serre.

d) \uline{Théorie des algèbres de Lie}: Comme sous-produit de recherches sur les
groupes algébriques en car. $p > 0$ , je trouve certains résultats délicats sur
les sous-algèbres de Borel ou de Cartan de certaines algèbres de Lie, notamment
sur les corps de base imparfaits (cf. [16, SGA 6, Exp. XIII et XIV]).

\uline{Références} : M. Demazure, J. Tits, J.P. Serre.

5. \uline{Géométrie Analytique} ([10], [11], [16 bis])

Mon influence sur la géométrie analytique est due moins aux résultats
nouveaux que j'ai pu y démontrer (la plupart contenus dans les réf. cit.), que
par les points de vue directement inspirés par la géométrie algébrique que j'ai
pu y introduire, et les nombreuses suggestions d'énoncés que j'ai pu y faire.
Un des plus anciens est le théorème de finitude de Grauert pour les morphismes
propres d'espaces analytiques, aboutissant à sa généralisation récente en un

\margin{199}
